\section{Билет 43. Первое начало термодинамики. Элементарная и полная работа. Графический смысл работы газа. Работа идеального газа при изобарном и изохорическом процессах (с выводом).}

\begin{definition}
\textbf{Первое начало термодинамики} является обобщением закона сохранения энергии на тепловые процессы. Оно гласит: количество теплоты $Q$, сообщаемое системе, расходуется на изменение её внутренней энергии $\Delta U$ и на совершение системой работы $A$ над внешними телами:
\begin{equation}
    Q = \Delta U + A. \label{eq:first_law_43}
\end{equation}
В дифференциальной форме для элементарных процессов: $\delta Q = dU + \delta A$. Знак $Q>0$ означает нагревание системы, $A>0$ --- совершение работы газом при расширении.
\end{definition}

\begin{definition}
\textbf{Элементарная работа} $\delta A$, совершаемая газом при квазистатическом расширении на малый объём $dV$, равна произведению давления $p$ на изменение объёма:
\begin{equation}
    \delta A = p \, dV. \label{eq:elem_work_43}
\end{equation}
\textbf{Полная работа} при конечном процессе перехода из состояния 1 в состояние 2 находится интегрированием:
\begin{equation}
    A = \int_{V_1}^{V_2} p(V) \, dV. \label{eq:total_work_43}
\end{equation}
\textbf{Графический смысл работы:} На термодинамической диаграмме $(p, V)$ работа численно равна площади фигуры, ограниченной кривой процесса, осью абсцисс ($V$) и вертикальными прямыми $V=V_1$, $V=V_2$.
\end{definition}

\begin{theorem}[Работа в изохорном процессе]
Изохорный процесс протекает при постоянном объёме: $V = \text{const}$, следовательно $dV = 0$. Подставляя это в \eqref{eq:elem_work_43}, получаем:
\begin{equation}
    \delta A = p \cdot 0 = 0 \quad \Rightarrow \quad A = 0. \label{eq:work_isochoric}
\end{equation}
Газ не совершает механической работы, так как отсутствуют перемещения границ системы. Вся подведённая теплота идёт на изменение внутренней энергии: $Q = \Delta U$.
\end{theorem}

\begin{theorem}[Работа в изобарном процессе]
Изобарный процесс протекает при постоянном давлении: $p = \text{const}$. Выносим $p$ за знак интеграла в \eqref{eq:total_work_43}:
\begin{equation}
    A = p \int_{V_1}^{V_2} dV = p (V_2 - V_1) = p \Delta V. \label{eq:work_isobaric_V}
\end{equation}
Используя уравнение состояния идеального газа $pV = \nu RT$, выразим работу через изменение температуры:
\begin{equation}
    p V_1 = \nu R T_1, \quad p V_2 = \nu R T_2 \quad \Rightarrow \quad p(V_2 - V_1) = \nu R (T_2 - T_1).
\end{equation}
Окончательно получаем формулу для изобарной работы:
\begin{equation}
    A = \nu R \Delta T. \label{eq:work_isobaric_T}
\end{equation}
Работа прямо пропорциональна изменению температуры и количеству вещества.
\end{theorem}

\begin{remark}
\textbf{Графическая иллюстрация:}
\begin{itemize}
    \item На $(p, V)$-диаграмме изохора --- вертикальная прямая. Площадь под ней равна нулю, что подтверждает $A=0$.
    \item Изобара --- горизонтальная прямая. Площадь под ней представляет собой прямоугольник со сторонами $p$ и $\Delta V$, площадь которого равна $p \Delta V$.
    \item При расширении ($\Delta V > 0$) работа газа положительна ($A>0$). При сжатии ($\Delta V < 0$) работа отрицательна ($A<0$), то есть внешние силы совершают работу над газом.
\end{itemize}
\end{remark}
